\begin{exercise}[No-arbitrage bounds for European calls and puts]
For \(S,K>0\) and \(t\in[0,T]\), prove the standard lower and upper bounds
for a European call \(C_t(S,T,K)\) and a European put \(P_t(S,T,K)\).
\end{exercise}

\begin{solution}
Let \(\tau=T-t\).  We use the comparison principle in absence of arbitrage:
if one terminal payoff is always dominated by another, then its time \(t\)
price is dominated as well; see \polyref{Prop.~1.2 and Cor.~1.1}.  The
terminal payoff of the call is
\[
  C_T=\pos{S_T-K}.
\]
First,
\[
  0\le \pos{S_T-K}\le S_T.
\]
The payoff \(S_T\) is obtained by holding one unit of the underlying from
\(t\) to \(T\), whose price is \(S\) at time \(t\).  Hence
\[
  0\le C_t(S,T,K)\le S.
\]
Moreover \(\pos{x}\ge x\) for every real number \(x\), so
\[
  \pos{S_T-K}\ge S_T-K.
\]
The payoff \(S_T-K\) is replicated at time \(t\) by buying one unit of the
underlying and borrowing \(Ke^{-r\tau}\) in the bank account.  Its time \(t\)
value is \(S-Ke^{-r\tau}\).  Therefore
\[
  C_t(S,T,K)\ge S-Ke^{-r\tau}.
\]
Combining this with \(C_t(S,T,K)\ge0\) gives
\[
  \boxed{\ \pos{S-Ke^{-r(T-t)}}\le C_t(S,T,K)\le S\ }.
\]

For the put, the terminal payoff is
\[
  P_T=\pos{K-S_T}.
\]
Since \(0\le \pos{K-S_T}\le K\), the comparison principle gives
\[
  0\le P_t(S,T,K)\le Ke^{-r\tau}.
\]
Also \(\pos{K-S_T}\ge K-S_T\).  The payoff \(K-S_T\) is replicated by
lending \(Ke^{-r\tau}\) at time \(t\) and short-selling one unit of the
underlying, so its time \(t\) value is \(Ke^{-r\tau}-S\).  Hence
\[
  P_t(S,T,K)\ge Ke^{-r\tau}-S.
\]
Together with non-negativity this yields
\[
  \boxed{\ \negpart{S-Ke^{-r(T-t)}}\le P_t(S,T,K)\le Ke^{-r(T-t)}\ }.
\]
These are exactly the bounds stated in \polyref{Prop.~1.4}.
\end{solution}
